\section{Noncommutative geometry and differential calculi}\label{sec:cohomology:noncommutative-geometry}

The differential calculi of this chapter start from functions that commute.
% claim-ok: own signpost from the cyclic homology section to this one
Noncommutative geometry drops that assumption.
The derivations themselves pay for it. The derivations of a
noncommutative ring $A$ fail to form an $A$-module
\cite[Abstract, p.~1]{Sardanashvily2007Dilemma}. % claim-ok: the cited sentence states the claim word for word

The $Q$-valued derivation of a noncommutative $K$-ring $A$ is a
$K$-module morphism $u\colon A\to Q$ with the Leibniz rule
$u(ab)=u(a)b+au(b)$ \cite[p.~1]{Sardanashvily2007Dilemma}. Several
definitions of a differential operator compete in this setting, and none
of them agrees with another \cite[Abstract, p.~1]{Sardanashvily2007Dilemma}.
The commutative side of the comparison is
Theorem~\ref{thm:pw:diffop-order}. % claim-ok: that theorem is the commutative order bound that Sardanashvily recalls on pp. 1--2

\begin{remark}[{\cite[Abstract, p.~1]{Sardanashvily2007Dilemma}}]
\label{rem:cohomology:ncg-dilemma}
In contrast with differential operators on modules over commutative and
graded commutative rings, there is no satisfactory notion of a
differential operator in noncommutative geometry
\cite[Abstract, p.~1]{Sardanashvily2007Dilemma}.
\end{remark}

The Chevalley--Eilenberg differential calculus still provides an
important example, and on smooth functions it is the de Rham complex
\cite[p.~1]{Sardanashvily2007Dilemma}. This is the smooth thread of
Section~\ref{sec:cohomology:computations} kept alive. % claim-ok: the cited page names the de Rham example and the section carries the comparison theorem

Ginzburg and Schedler give one answer on the operator side
\cite[Abstract, p.~673]{GinzburgSchedler2010BV}. They replace
the derivations of a commutative algebra by the bimodule
$\mathbb{D}\mathrm{er}(A)$ of \emph{double derivations} of an associative
algebra $A$. Their operators act on a Fock space $\mathcal{F}(A)$ built
from $A$ \cite[Abstract, p.~673]{GinzburgSchedler2010BV}.

\begin{remark}[{\cite[Abstract, p.~673]{GinzburgSchedler2010BV}}]
\label{rem:cohomology:ncg-fock}
The algebra $\mathcal{D}(\mathcal{F}(A))$ of Ginzburg and Schedler is
filtered, and its associated graded algebra is commutative in a wheeled
sense. For a smooth algebra $A$ it satisfies
$\operatorname{gr}\mathcal{D}(\mathcal{F}(A))\cong
\mathcal{F}(T^{A}(\mathbb{D}\mathrm{er}(A)))$. A Ricci-flat, torsion-free
bimodule connection on $\mathbb{D}\mathrm{er}(A)$ gives a second-order
wheeled differential operator, a noncommutative analogue of the
Koszul--Batalin--Vilkovisky delta operator
\cite[Abstract, p.~673]{GinzburgSchedler2010BV}.
\end{remark}

The derivation therefore survives in noncommutative geometry three times
over. As an object of study it loses its module structure
\cite[Abstract, p.~1]{Sardanashvily2007Dilemma}. As a double derivation
it becomes the building block of a new calculus
\cite[Abstract, p.~673]{GinzburgSchedler2010BV}. On smooth functions it
also keeps the de Rham calculus of the manifold
\cite[p.~1]{Sardanashvily2007Dilemma}. % claim-ok: closing summary of the two cited remarks above
